% Generated draft: requires mathematical and source review.
\documentclass[11pt]{article}
\usepackage[margin=25mm]{geometry}
\usepackage{fontspec}
\setmainfont{Latin Modern Roman}
\usepackage{amsmath,amssymb,mathtools,unicode-math}
\setmathfont{Latin Modern Math}
\usepackage{enumitem,needspace,xurl,hyperref}
\setlength{\parindent}{0pt}
\setlength{\parskip}{5pt}
\setlength{\emergencystretch}{3em}
\newcommand{\N}{\mathbb{N}}
\newcommand{\Z}{\mathbb{Z}}
\newcommand{\Q}{\mathbb{Q}}
\newcommand{\R}{\mathbb{R}}
\newcommand{\C}{\mathbb{C}}
\newcommand{\F}{\mathbb{F}}
\newcommand{\A}{\mathbb{A}}
\newcommand{\PP}{\mathbb{P}}
\newcommand{\E}{\mathbb{E}}
\newcommand{\eps}{\varepsilon}
\newcommand{\dd}{\,\mathrm{d}}
\begin{document}
\section*{UnsolvedMath ID: 30003081}
\textbf{Source problem code:} catalogue entry
\subsection{Definitions and mathematical statement}
For notation only, let \(k\) be a field, let \(X\) be a smooth \(k\)-variety, and let \(D\subset X\) be a reduced divisor with ideal sheaf \(\mathcal I_D\). The sheaf of logarithmic vector fields along \(D\) is
\[
\mathcal T_X(-\log D)
=
\{\theta\in\mathcal{D}er_k(\mathcal O_X,\mathcal O_X):
\theta(\mathcal I_D)\subseteq\mathcal I_D\}.
\]
These choices explain the terminology; the supplied record does not specify the base field, characteristic, varieties, or divisors required in the intended extension.

A standard possible meaning of an unexpected hypersurface can be described conditionally as follows. For a finite reduced set \(Z\subset\PP^n_k\), integers \(d,m\geq1\), and a point \(P\notin Z\), put
\[
V_d(Z)=H^0(\PP^n_k,\mathcal I_Z(d)),\qquad
V_{d,m}(Z,P)=H^0(\PP^n_k,(\mathcal I_Z\cap\mathcal I_P^m)(d)).
\]
Under the convention that a general point of multiplicity \(m\) is expected to impose \(\binom{n+m-1}{n}\) independent conditions, unexpectedness means
\[
\dim_k V_{d,m}(Z,P)>
\max\left\{0,\dim_k V_d(Z)-\binom{n+m-1}{n}\right\}
\]
for general \(P\). Here general means throughout a nonempty Zariski open subset of the parameter space of points, interpreted geometrically if necessary. This is an explanatory convention, not a definition supplied by the record; whether it is the intended notion, and which field hypotheses are required, remain unresolved.

\textbf{Source problem.} Develop a higher-dimensional analogue of the logarithmic-vector-field exact sequence used in the plane-curve characterization, and determine whether it gives corresponding characterizations of unexpected hypersurfaces. The original wording asks whether the theorem characterizing unexpected curves through multiplicity indices extends to higher dimensions using the same methods. For projective ambient spaces, higher than the plane would mean \(\PP^n\) with \(n\geq3\), but the record does not specify the range or hypotheses of the proposed extension.

The plane theorem, the exact sequence and its maps, the multiplicity indices, and their proposed higher-dimensional replacements are absent from the supplied record. Consequently, this is a statement of the requested research direction with explicit unresolved qualifications, rather than a uniquely specified quantified conjecture. No equivalence between the two supplied formulations, or higher-dimensional characterization, is asserted.
\subsection{Short English statement}
The problem asks whether the logarithmic-vector-field methods behind a characterization of unexpected plane curves extend to unexpected hypersurfaces in higher dimensions. The original formulation emphasizes multiplicity indices. The exact plane theorem, sequence, hypotheses, and intended higher-dimensional invariants are missing, so a precise theorem-shaped conjecture requires further source clarification. The supplied open label is a dated curation assessment.
\subsection{Sources}
\begin{itemize}
\item Supplied local UnsolvedMath registry record: permanent ID \(30003081\), title \emph{Higher-Dimensional Theory of Unexpected Curves}, source problem code \texttt{catalogue entry}, in the Oberwolfach Reports Open Problems collection. The record gives proposed year 2016 and names no individual proposer.
\item \emph{Mini-Workshop: Arrangements of Subvarieties, and their Applications in Algebraic Geometry} (2016), Oberwolfach Reports, pp.~653--702, DOI 10.4171/owr/2016/14. Supplied report reference and literature-source link: \url{https://doi.org/10.4171/owr/2016/14}. The background describes this report as documenting workshop discussions and open problems; no precise problem page or theorem number is supplied.
\item Schenck--Terao--Yoshinaga, supplied literature-source link: \url{https://arxiv.org/abs/1407.3237}. The record supplies no title or specific result locator for this reference.
\item Literature assessment dated exactly \texttt{2026-08-21}: the record labels the problem open and states that the requested higher-dimensional analogue is open while the cited methods establish the plane-curve case only. This is supplied dated curation, not an independently verified assessment; the background explicitly requests independent verification of time-sensitive status and cited claims.
\item The registry labels the statement \texttt{corrected\_verified} and reports that the official report passage was checked and rewritten while preserving its mathematical content. This records the supplied verification claim; no independent source visit or verification is asserted here.
\end{itemize}
\end{document}
